\documentclass[../../../include/open-logic-section]{subfiles}

\begin{document}

\olfileid{sth}{choice}{banach}
\olsection{مفارقة باناخ--تارسكي}

وقد نلتمس للاختيار مسوغًا خارجيًا، كما التمس بولوس للاستبدال؛ وانظر
\olref[replacement][extrinsic]{sec}. فإن قبول الاختيار يثمر أمورًا
نافعة جمة: مقارنة كل عددين كارديناليين، وترتيب كل مجموعة ترتيبًا جيدًا،
وجعل الأعداد الكاردينالية ضربًا مخصوصًا من الأعداد الترتيبية، وغير ذلك.

غير أن للاختيار، فيما يقال، نتائج مستقبحة. وأكثر هذا القول راجع إلى
نتيجة~\cite{BanachTarski1924}.

\begin{thm}[مفارقة باناخ--تارسكي (في $\ZFC$)]
يمكن أن تفكك أية كرة قطعًا متناهية العدد، ثم تعاد هذه القطع، بالدوران
والإزاحة، فتؤلف كرتين مماثلتين للكرة الأولى.
\end{thm}
\noindent
وهذا في بادئ النظر عجب؛ فإن حجم الكرتين ضعف حجم الكرة التي بدأنا بها،
مع أن الحركات الصلبة من دوران وإزاحة لا تبدل الحجم. فكأن مبرهنة
باناخ--تارسكي تخرج لنا مادة من العدم.

ثم الأمر أعجب من ذلك.\footnote{انظر~\citet[Theorem 3.12]{Wagon2016}.}
فبحجة من جنسها يمكن تقطيع حبة بازلاء قطعًا متناهية، ثم تدوير القطع
وإزاحتها وإعادة جمعها حتى تصير جسمًا في شكل ساعة بيغ بن وحجمها.

ولا يثبت شيء من ذلك في~$\ZF$ وحدها.\footnote{مع أنه يمكن إثبات
باناخ--تارسكي بمبادئ أضعف تمامًا من الاختيار؛ انظر
\citet[303]{Wagon2016}.} فنحن بين أمرين: أن نرد الاختيار، أو أن نألف
هذه «المفارقة».

والذي نختاره أن نألفها؛ بل لا نعدها مفارقة عظيمة. وليس يظهر لنا وجه
كونها أشد مفارقة أو أخف من الأحكام الآتية.\footnote{يسوق كل من
\citet[276--7]{Potter2004} و\citet[16]{Weston2003} و
\citet[31, 308--9]{Wagon2016} ملاحظات مشابهة، مستعملًا أمثلة أخرى.}
\begin{enumerate}
	\item عدد النقاط في الفترة~$(\OLMathNumeral{0},\OLMathNumeral{1})$ مساو لعددها في~$\Real$.
	\\\emph{البرهان}: انظر في~$\tan(\pi(\OLId{latin-ordinary}{r}-\nicefrac{\OLMathNumeral{1}}{\OLMathNumeral{2}}))$.
	\item عدد النقاط في خط مساو لعددها في مربع.
	\\انظر~\olref[his][set][pathology]{sec} و\olref[his][set][cantorplane]{sec}.
	\item توجد منحنيات تملأ الفضاء.
	\\انظر~\olref[his][set][pathology]{sec} و\olref[his][set][hilbertcurve]{sec}.
\end{enumerate}
ولا يفتقر واحد من هذه الأحكام الثلاثة إلى الاختيار، وقد صرنا نعدها من
محاسن الرياضيات المدهشة. فلعل مفارقة باناخ--تارسكي جديرة بالحكم نفسه.

نعم، لا بد هنا من تنبيه تقني لا يحتاج إلا إلى روية: الحركات الصلبة
تحفظ الحجم، فلا يمكن أن تكون القطع الخمس\footnote{صغنا المفارقة
بعبارة «عدد منته من القطع». والواقع أن
\citet{Robinson1947} أثبت إمكان إنجاز التفكيك بـ\emph{خمس} قطع (ولا
يمكن بأقل منها). ولبرهان ذلك، انظر~\citet[pp.~66--7]{Wagon2016}.}
التي نفكك إليها الكرة قابلة للقياس كلها. وعلى التقريب، لا يصح أن نسند
إلى كل قطعة منها حجمًا. وإنما ينبغي أن تتصورها مجموعات متناثرة من
النقاط، لا متناهية ولا قابلة للتصوير. وقد يبدو وجود مثل هذه المجموعات
غريبًا، ولكنه يلزم فيما يظهر عن الأمر الكامن في~\stagesacc{} بتكوين كل
مجموعة ممكنة من المجموعات المتاحة من قبل.

فإن لم يكفك هذا كله، فهذه حجة خارجية أخيرة لقبول مفارقة
باناخ--تارسكي: إنها تنتج من فورها أحسن نكتة رياضية.
\begin{enumerate}
	\item[] \emph{سؤال}. ما إعادة ترتيب حروف «باناخ--تارسكي»؟
	\item[] \emph{جواب}. «باناخ--تارسكي باناخ--تارسكي».
\end{enumerate}

\end{document}
